\section{Decision shortcuts and the encoding obstacle}
Reducing the root modulus to two does not remove the current algorithm's internal precision cost. Its recurrence has the form
\[
 p_t=p+8\lceil3M/u\rceil\delta+8M(a+1).
\]
Even $p=1$ leaves the $\Omega(M\delta/u)$ term. Substituting the existing precision certificate into the mode estimate still gives the same $O(M/\sqrt u)$ exponent. This does not exclude a different parity-counting algorithm; it explains why parity alone does not improve this one.

There is also a precise limitation on lowering polynomial degree while preserving witnesses.
\begin{proposition}[Degree needed to preserve an odd solution count]\label{prop:odd}
For polynomials $P_1,\ldots,P_a$ over $\F_2$ on $n$ Boolean variables, if $\sum_i\deg P_i<n$, their number of common zeros is even. Consequently, a same-variable solution-set-preserving encoding of a uniquely satisfiable formula cannot always have total degree $o(n)$.
\end{proposition}
\begin{proof}
On the Boolean cube the common-zero indicator is
$Q(x)=\prod_i(1+P_i(x))$. Its degree is less than $n$. Every monomial in $Q$ omits some variable, so summing it over all Boolean assignments gives zero in $\F_2$. Thus $\sum_xQ(x)=0$. A unique common zero would instead give one.
\end{proof}
This statement does not constrain arbitrary decision reductions or reductions that intentionally change the witnesses.

The incidence encoding creates a different obstacle. If different rows use disjoint incidence variables and each has $r$ distinct cubic monomials, the product of $a$ row factors $(1+P_i)$ has at least $r^a$ distinct incidence monomials: independently choose one cubic term in each row. The rows' disjoint indeterminates prevent collisions or cancellation among these products. With $a=\Theta(n)$ and $r=\log^{\Theta(1)}n$, injectively indexing these monomials needs $\Omega(n\log\log n)$ bits.

This is not restricted to an arbitrary unrelated polynomial example. The bundled note \texttt{research\_decision/findings.md}, Section 2, gives a complete probabilistic construction from genuine 3-CNF violation polynomials. For $n\ge\lceil e^{10}\rceil$, choose $m=n\lceil\ln n\rceil$ clauses on distinct variable triples and $6n$ parity rows, including each clause independently with probability $4\ln(2m)/n$. A union bound over violation sets and assignments, together with row-size tail bounds, shows that some matrix preserves the exact solution set and has $\Theta(\log^2 n)$ distinct cubic terms per row. Its separated-incidence product has $2^{\Omega(n\log\log n)}$ nonzero monomials. The construction is an existence argument, not a proposed randomized solver or a description of the paper's deterministic greedy output.

Identifying incidence copies with their original Boolean variables may create very large cancellations and leaves at most $2^n$ multilinear supports. An algorithm that performs this quotient implicitly while retaining small local descriptions and low degree would evade the incidence-indexing obstruction. Such an algorithm has not been established here.

Other black-box retunings do not fix the problem. The supplied sparsification certificate, followed only by direct fixed-density contraction, optimizes for 3-SAT to exponent
\[
 O\!\left(\frac{n(\log\log n)^{2/3}}{(\log n)^{1/6}}\right),
\]
which is worse than \eqref{eq:baseline}. Branching on $t$ variables and paying the same black-box cost on $n-t$ variables gives exponent $t+Cf(n-t)$, where $f(n)=n\log\log n/\sqrt{\log n}$. Since $f'(n)\to0$, branching on $t\le n/2$ eventually increases that certificate; $t>n/2$ already costs exponentially more than $2^{Cf(n)}$. These conclusions concern those specific certificates, not all sparsification or branching algorithms.

% Research note; no unconditional improvement is asserted.
\section{A conditional degree-preserving compression theorem}
\label{sec:conditional-degree-kernel}
All logarithms in this section are base two, except in analytic partition estimates.
Let $s$ and $K$ be positive integers, with $s$ a size parameter, $\lambda=\lceil\log_2 s\rceil$, and suppose
$s/\lambda\le K\le s$. Constants multiplying these inequalities can be absorbed
by changing $s$ and the constants below. Put
\[
 H=\left\lceil\sqrt{sK/\lambda}\right\rceil,
 \quad u_d=\lfloor\varepsilon\lambda/2^d\rfloor,
 \quad a_d=2\sum_{j=0}^d u_j,
 \quad L_d=2^{a_d},\qquad \varepsilon=(128\cdot1024)^{-1}.
\]
The parameter $K$ is fixed globally; it need not decrease in a child.

\paragraph{Explicit hypothesis, not an established construction.}
Assume a class of encoded contraction objects admits the following uniform
recursive interface. Every object has $M\le Cs8^{-d}$ Boolean wires at depth $d$,
maximum degree at most $\lceil s^{9/10}\rceil$, and entries over the ring used in
Appendix D. Root precision and nilpotent truncation order are $O(s)$.
Here $\mathcal L$ denotes the total root description length. The graph and
immutable entry descriptions stored along any root-to-leaf path have total
length at most $(s+\mathcal L)^{O(1)}$, with a fixed uniform exponent, and all
entry evaluations, at their creation precisions, take
$\operatorname{poly}(s+\mathcal L)2^{O(s^{9/10})}$ time and polynomial space.
Each active object, $M>H$, must have all the following properties.
\begin{enumerate}
\item There is an exactly specified deformation $F(T)$ whose value at zero
is its contraction and whose degree is at most $K(2L_d-1)$. This is the key
new degree bound; the published bit-wire deformation only guarantees
$M(2L_d-1)$. For $S_d=\{-L_d,\ldots,L_d-1\}$ and
$\alpha_d=2^{a_d+1}$, all $F(t+\alpha_d z)$ are integral over the coefficient
ring. Their individual local factors and evaluations must be constructible
at polynomial bit cost, apart from the stated entry-evaluation cost.
\item At every retained node $t\ne0$ with $i=v_2(t)<u_d$, put
$D_t=L_d/2^i$. At any residue-field interpolation value of $z$, the evaluation
has an exact expansion modulo its working precision into mode monomials.
There are at most $cM/D_t$ modes of each positive integer weight $w$;
a mode monomial of weighted degree $b$ has coefficient divisible by $2^b$.
Modes and feasible exponent tuples have polynomial descriptions and can be
enumerated with polynomial delay and space. Merely showing $cN/D_t$ modes
for an uncompressed incidence count $N$ is not sufficient.
\item For each retained monomial, its contribution factors as the product
of a constant number of child contractions, each with at most $M/8$ wires.
The children remain in this class, with the same global degree parameter $K$,
maximum degree at most $\lceil s^{9/10}\rceil$, and the description and entry
cost bounds above. Entrywise canonical lifts must retain their creation-time
representatives, as in Appendix D. In particular, this condition prohibits
silently replacing a high-arity group tensor by an equality tree.
\item All local interpolation, packing, mode processing, rational recovery,
and description work per generated term is polynomial in $s+\mathcal L$;
the number of interpolation nodes, evaluation points, and groups per call
is polynomial in $s$ alone, independently of $\mathcal L$. Leaves can be summed over their Boolean wires directly, with the
stated entry costs. The coefficient-ring arithmetic uses polynomial-bit
coordinates and polynomial space at the precisions derived below.
\end{enumerate}

\paragraph{Conditional conclusion.}
Under this interface, contraction takes
\[
 \operatorname{poly}(s+\mathcal L)\,
 2^{O(\sqrt{sK/\log s})}
\]
bit operations and polynomial space. This is an implication from a new
recursive interface, not a proof that the interface exists for 3SAT.

\paragraph{Proof of the implication.}
Use the published tapered Hermite lemma with its degree parameter $M$
replaced by $K$. Thus
\[
 h=\lceil3K/u_d\rceil,\quad q_t=h\max(0,u_d-i),\quad
 p_t=p+8hD_t+8K(a_d+1).
\]
Here $K\ge u_d$ for all sufficiently large $s$. There are enough jets,
because $\sum q_t\ge L_dhu_d\ge3KL_d>K(2L_d-1)$; also $h\le4K/u_d$.
The proof of the source lemma then gives exactly the same integral recovery.
One may recover each jet by polynomially many residue-distinct evaluations
because the degree is $O(KL_d)=s^{O(1)}$.

Since $H\ge s/\lambda$, active depth is $O(\log\lambda)$; more precisely,
$8^d<Cs/H\le C\lambda$, so $u_d=\Omega_C(\lambda^{2/3})$.
Also $a_d\le4\varepsilon\lambda$, and the original packing inequalities
remain true: their right sides are $O(s^{0.902})=o(H)$ or smaller.
Summing previous precision increments gives
\[
 p\le s\lambda^{O(1)}2^{a_{d-1}}.
\]
For a retained node,
$D_t\ge2^{a_d-u_d}=2^{a_{d-1}+u_d}$.
As $K\ge s/\lambda$ and $2^{u_d}$ dominates every fixed power of $\lambda$,
the inherited precision and the $8K(a_d+1)$ term are absorbed, yielding
\[
 p_t=O(KD_t/u_d).
\]
All precisions remain polynomial in $s$ because $a_d=O(\log s)$.

The generating-function estimate from Appendix D, now with at most $cM/D_t$
variables of every weight, gives
\[
 \#\{\beta:\sum_Zw_Z\beta_Z<p_t\}
 \le \exp\bigl(O(\sqrt{Mp_t/D_t})\bigr)
 =\exp\bigl(O(\sqrt{MK/u_d})\bigr).
\]
At depth $d$ this exponent is at most
\[
 O\!\left(\sqrt{sK/\lambda}\,2^{-d}\right),
\]
because $M\le Cs8^{-d}$ and $u_d=\Theta_C(\lambda2^{-d})$.
Summing along paths, and accounting for the polynomial branching factors,
gives $\exp(O(\sqrt{sK/\lambda})+O(\lambda\log\lambda))$ calls.
The latter additive term is negligible since $\sqrt{sK/\lambda}\ge s/\lambda$.
Leaf work is bounded by $\operatorname{poly}(s+\mathcal L)
\exp(O(H)+O(s^{9/10}))$, with the same final exponent.
Depth-first accumulation and immutable polynomial descriptions give
polynomial space. This proves the conditional implication.

\paragraph{Consequent target for the sparse-query reduction.}
The source reduction has compressed wire budget $s=\Theta(n\nu)$ and retained
incidence degree $K=\Theta(n)$, where $\nu=\lceil\log_2\lceil\log_2 n\rceil\rceil$.
If its encoded objects and every recursive child admitted the above
interface, the source's other overheads would still be absorbed, giving
\[
 (n+m)^{O(1)}2^{O(n\sqrt{\log\log n/\log n})}
\]
for arbitrary-density exact counting and hence decision. This improves the
published exponent by a factor $\sqrt{\log\log n}$. It does not improve the
published fixed-density $O(n/\sqrt{\log n})$ exponent.
The retained incidence degree is not, by itself, a degree bound on the
multilinear extension of a binary subset index. No argument establishing
that implication or recursive closure has been found. Already the coefficient
table of the degree-one incidence polynomial $z_{2^b}$ on singleton indices,
encoded by $b$ bits in ordinary order, is the Boolean AND of all $b$ index
bits, of degree $b$. This elementary example rules out a generic degree
transfer based solely on the incidence degree.

\paragraph{A local candidate and its obstruction.}
For a categorical wire indexing $d$-element subsets of an $N$-element set,
consider $K_{S,T}(T_0)=\det W(T_0)[S,T]$, the $d$th exterior power of a base
matrix deformation. If $W(0)=I$, then $K_{S,T}(0)=\delta_{S,T}$; its degree is
at most $d\deg W$, and a block-diagonal $W$ induces block factorization after
group occupancies and determinant signs are tracked. Integrality of $W$
implies integrality of these determinants. Thus this is a genuine local
low-degree kernel, not yet the required recursive contraction algorithm.
Its ordinary Cauchy mode construction uses $N/D$, not
$\log\binom Nd/D$, modes per weight. Exterior antisymmetry can force a
rank-one mode to occur at most once in one determinant, but distinct-part
partition counts still have square-root entropy when the particle bound
is inactive. The missing compressed mode multiplicity and child-class
closure are substantive new lemmas.



\section{An exact sufficient compression lemma}\label{sec:conditional}
The following hypothesis isolates a sufficient missing theorem. It is deliberately stated with all interface costs rather than just a wire count.

\begin{hypothesis}[Uniform streaming compression]\label{hyp:compression}
Fix a function $W$ with $n\le W(n)\le n(\log n)^c$ for a fixed constant $c$. There is a deterministic polynomial-space procedure that, for every 3-CNF $F$, streams at most
\[
 \poly(n,m)2^{O(W(n)/\sqrt{\log W(n)})}
\]
Boolean tensor networks in the same total time. Uniform fixed constants guarantee the following for each network:
\begin{enumerate}
\item It is bipartite, has at most $CW(n)$ wires, at most polynomially many vertices in $W(n)$, and maximum degree at most $\lceil W(n)^{9/10}\rceil$.
\item Entries belong to $(\Z/2^{n+1}\Z)[V]/(V^\ell)$, where $1\le\ell\le CW(n)$. Write $\mathcal L$ for the total graph and evaluator-description length, which is polynomial in $n+m$.
\item Every entry at every Boolean port assignment is computable in time $(W(n)+\mathcal L)^b2^{AW(n)^{9/10}}$ and space $(W(n)+\mathcal L)^b$, with fixed constants $A,b$ independent of $F,n,m$. The immutable description may include the original input encoding; this adds no wires.
\item A specified polynomial-time linear functional of each contraction gives a residue; the sum of these residues is $\#F$ modulo $2^{n+1}$.
\end{enumerate}
\end{hypothesis}

\begin{proposition}[Conditional consequence]\label{prop:compression}
Assuming the paper's contraction theorem and Hypothesis~\ref{hyp:compression}, exact counting and decision take time
\[
 \poly(n,m)2^{O(W(n)/\sqrt{\log W(n)})}
\]
and polynomial space. In particular, $W(n)=O(n)$ would improve the unrestricted-density exponent to $O(n/\sqrt{\log n})$.
\end{proposition}
\begin{proof}
Apply the contraction theorem with $s=\lceil W(n)\rceil$ to one network at a time. Its initial precision $n+1$ is $O(s)$, and every other hypothesis is explicitly supplied. Multiplying the per-call bound by the number of networks changes only the fixed constant in the exponential rate. Apply the given functionals and accumulate modulo $2^{n+1}$. The true count lies between zero and $2^n$, so this residue determines the count exactly. Streaming preserves polynomial space.
\end{proof}

Any $W(n)=o(n\log\log n)$ in the stated range already improves the general exponent, because $\log W(n)=\Theta(\log n)$. The hard missing part is Hypothesis~\ref{hyp:compression}, not Proposition~\ref{prop:compression}. Merging the entire formula into one $n$-port tensor does not prove it: the degree cap fails, and evaluating that tensor after a purported degree reduction may conceal the original hard computation.

A concrete place to search is the paper's older original-variable interface before incidence expansion. It already uses $n$ Boolean variables and total input degree $O(n)$ but has the weaker loglog-squared analysis. A recursive replacement would need to prove that transformed factors retain controlled local degree, that shared measurement variables do not restore the incidence-code size, that local evaluation stays subdominant, and that interpolation precision remains controlled at every depth. None of these closure properties follows just from having $n$ original variables.

